% sections/conclusion.tex
\chapter{Заключение}

В ходе работы были изучены и реализованы два классических алгоритма:

\begin{itemize}
  \item \textbf{Сортировка слиянием}~--- рекурсивный алгоритм с гарантированной
        сложностью $O(n \log n)$ (формула~\eqref{eq:merge-complexity}),
        пригодный для внешней сортировки больших объёмов данных.
  \item \textbf{Алгоритм Дейкстры}~--- эффективный жадный алгоритм поиска
        кратчайших путей (формула~\eqref{eq:dijkstra-complexity}),
        широко применяемый в сетевой маршрутизации и навигации.
\end{itemize}

Исходный код на Python (приложение~\ref{app:merge-sort-full}) и C++
(приложение~\ref{app:dijkstra-full}) полностью соответствует описанным
алгоритмам и протестирован на примерах из данной работы.

Результаты сравнения алгоритмов сортировки (таблица~\ref{tab:sort-comparison})
демонстрируют преимущества сортировки слиянием в плане стабильности
временной сложности по сравнению с алгоритмами, имеющими квадратичный
худший случай.
